% Solutions/hints for ch03 exercises (assembled into Appendix C)

\begin{solution}{exr:ch03-trace}
With $c(B,E)=4$ the second step pushes $(6,E)$ instead of $(9,E)$, and the rest of the trace changes accordingly. The pops are $(0,A)$, $(2,B)$, $(3,C)$, $(5,C)$ stale, $(5,D)$, $(6,D)$ stale, $(6,E)$, $(8,G)$, $(9,F)$, $(10,F)$ stale and $(13,F)$ stale: 11 pops and 4 stale entries. When $D$ is expanded, $E$ already has the tentative cost $6<5+3$, so $D$ pushes only $(13,F)$; when $E$ is expanded it pushes $(10,F)$ and $(8,G)$; $G$ then improves $F$ to 9. The final distances are $\dist(A,B)=2$, $\dist(A,C)=3$, $\dist(A,D)=5$, $\dist(A,E)=6$, $\dist(A,G)=8$ and $\dist(A,F)=9$, and the shortest path to $G$ is $A,B,E,G$ of cost 8. The self-test of \code{ch03\_dijkstra.py} contains this variant.
\end{solution}

\begin{solution}{exr:ch03-early}
$F$ is pushed for the first time in step 5 with cost 13 and parent $D$, so a search that stops at discovery reports $\dist(A,F)=13$ and the path $A,B,C,D,F$, whereas the true distance is $11$ along $A,B,C,D,E,G,F$. For $G$ the first push, in step 7, already carries the optimal cost 10, but only because no cheaper route to $G$ exists; the algorithm has no way of knowing that at the time. The proof of \cref{thm:ch03-invariant} compares the key of $u$ with the keys still in the heap \emph{because the heap returned $u$ as its minimum}. A vertex that has merely been pushed has not been returned by the heap, so nothing relates its tentative cost to the keys of the other open vertices; the only guarantee it has is $\gcost(v)\ge\dist(s,v)$, which is the wrong direction.
\end{solution}

\begin{solution}{exr:ch03-shift}
Take the graph of \cref{fig:ch03-negative-edge} and add $K=2$ to every edge: $S\to A$ costs 3, $S\to B$ costs 4, $A\to T$ costs 3 and $B\to A$ costs 0. In the shifted graph $S,A,T$ costs 6 and $S,B,A,T$ costs 7, so the algorithm returns $S,A,T$; in the original graph these paths cost 2 and 1, and $S,B,A,T$ is optimal. A path with $k$ edges gains exactly $kK$, so the shift penalises paths with more edges. It is harmless only when every path under comparison has the same number of edges, which is the case, for example, between two layers of the time-expanded graph of \cref{ch:ch02}, where every path from time 0 to time $T$ has exactly $T$ edges.
\end{solution}

\begin{solution}{exr:ch03-backward}
(a) The backward run from $G$ settles, in this order, $\hcost^*(G)=0$, $\hcost^*(F)=1$, $\hcost^*(E)=2$, $\hcost^*(D)=5$, $\hcost^*(C)=7$, $\hcost^*(B)=8$, $\hcost^*(A)=10$. Consistency holds on every edge, with equality on the edges of the backward tree: for instance $\hcost^*(D)=5\le c(D,F)+\hcost^*(F)=8+1$, $\hcost^*(A)=10\le c(A,C)+\hcost^*(C)=5+7$ and $\hcost^*(A)=10=c(A,B)+\hcost^*(B)=2+8$. Note that $\hcost^*(A)=\dist(A,G)=10$ agrees with \cref{tab:ch03-trace}. (b) A move into a hatched cell costs 3 but a move out of it costs 1, so the cost of a path depends on its direction, and $\dist(T,v)\neq\dist(v,T)$ exactly when the cheapest route between $v$ and $T$ enters a different number of hatched cells in the two directions. Running \code{dijkstra(grid, T)} and \code{backward\_dijkstra\_heuristic(grid, T)} shows that the two tables differ on the following cells (row, column): $(1,1)$: $10$ forward against $8$ backward, $(1,2)$: $13$ forward against $11$ backward, $(1,3)$: $10$ forward against $8$ backward, $(2,1)$: $9$ forward against $7$ backward, $(2,3)$: $7$ forward against $5$ backward, $(3,1)$: $6$ forward against $4$ backward, $(3,2)$: $5$ forward against $3$ backward, $(3,3)$: $4$ forward against $2$ backward. (c) The constraint removes the single state $((4,2),6)$ from the space-time graph. Every path that respects it is still a space-time path, and dropping the time index turns it into a walk in the grid of no larger cost, so $\hcost^*(v)=\dist(v,T)$ remains a lower bound (item~3 of \cref{thm:ch03-true-distance}). It is no longer exact: with unit costs the optimal path of \cref{ex:ch03-grid} reaches $(4,2)$ at time step 6, precisely when the cell is forbidden, so the drone must wait one step or take a detour and, with a wait action of cost 1 as in the unit-cost setting of \cref{ch:ch02}, the true constrained cost from $S$ is 9, while $\hcost^*(S)=8$.
\end{solution}

\begin{solution}{exr:ch03-bidirectional}
(a) Take the undirected graph $S$--$a$ of cost 1, $a$--$b$ of cost 28, $b$--$T$ of cost 1, and $S$--$u$ of cost 16, $u$--$T$ of cost 16, and alternate the two sides. Forward settles $S$ (relaxing $a$ to 1 and $u$ to 16), backward settles $T$ (relaxing $b$ to 1 and $u$ to 16), forward settles $a$, backward settles $b$; then $\text{\Open}_f=\{u{:}16,b{:}29\}$ and $\text{\Open}_b=\{u{:}16,a{:}29\}$, and the next vertex settled by both sides is $u$, with $\gcost_f(u)+\gcost_b(u)=32$. But $\dist(S,T)=30$ along $S,a,b,T$, so the naive rule returns a path 2 units too expensive. The rule with $\mu$ does better: when the forward expansion of $a$ scanned the edge $(a,b)$, $\gcost_b(b)=1$ was already finite, so $\mu=1+28+1=30$ was recorded, and \cref{eq:ch03-bidir-stop} fires at $16+16\ge30$ --- before $u$ is settled twice --- returning $S,a,b,T$.
(b) Let $\pi$ be an $S$--$T$ path with $c(\pi)<\mu$ when the rule fires. If some vertex of $\pi$ is still open on the forward side and a later one is still open on the backward side, let $u$ be the first vertex of $\pi$ not settled forward and $w$ the last one not settled backward, so $u$ comes no later than $w$. The predecessor of $u$ on $\pi$ is settled and expanded forward, so $u$ carries the key $\gcost_f(u)=\dist(S,u)$ in $\text{\Open}_f$ (step~(ii) of the proof of \cref{thm:ch03-invariant}), and symmetrically $\gcost_b(w)=\dist(w,T)$ in $\text{\Open}_b$. Since the prefix of $\pi$ up to $u$ costs at least $\dist(S,u)$, the suffix from $w$ at least $\dist(w,T)$ and the piece between them at least 0, $c(\pi)\ge\min_{\text{\Open}_f}\gcost_f+\min_{\text{\Open}_b}\gcost_b\ge\mu$, a contradiction. Otherwise the forward-settled prefix and the backward-settled suffix of $\pi$ touch, so $\pi$ has an edge $(x,y)$ with $x$ settled forward and $y$ settled backward. Whichever of the two was settled later scanned $(x,y)$ during its own expansion, when $\gcost_f(x)=\dist(S,x)$ and $\gcost_b(y)=\dist(y,T)$ were both final, so $\mu\le\gcost_f(x)+c(x,y)+\gcost_b(y)\le c(\pi)$, again a contradiction. Hence no path is cheaper than $\mu$.
(c) Report the two settled-cell counts side by side for each grid. The disc argument predicts a factor of about one half on open maps; the measured factor is worse whenever obstacles funnel both waves through the same corridor, since the two discs then overlap long before they meet.
\end{solution}
